\chapter*{摘要}\addcontentsline{toc}{chapter}{摘要}
本项目围绕机器人运动控制中的轨迹、测量、接触和支撑问题，逐步完成ROS2仿真通信、Panda关节轨迹跟踪、惯性测量单元（IMU）滤波与实体采集、末端导纳控制、倒立摆平衡、平面双足行走及三维人形全身控制。技术路线从容易隔离问题的简化工况开始，逐步恢复重力、碰撞、浮动基座和有限电机力矩，并用失败对照检查方案的适用范围。

机械臂部分，5秒五次轨迹经过驱动增益与参考速度调整，在独立幅度验证中的平均运动误差由0.864度降至0.082度。实体IMU在Jetson上的新测量约25赫兹；提高到200赫兹发布时约87.52\%消息复用旧测量，说明发布频率与信息更新率应分别评价。接触控制从平均力改善推进到接近减速、目标力缓升和重力补偿。选用刚性接触条件的1毫秒原始峰值为11.11牛，末段均力约5.007牛；物理步长减半时原始峰值升至20.43牛，仍存在数值敏感性。

双足部分，先在二维线性倒立摆中验证落脚调整与因果延迟预测，再将连续摆腿、双支撑和逆动力学二次规划接入MuJoCo多关节动力学。后期平面模型完成两组二十步行走，前进约2.323米；94个唯一条件通过79个，较长延迟仍导致明显失败。三维G1模型保留自由浮动基座，以29个受限电机力矩完成站立、双侧单脚支撑和慢速交替迈步。正式52个条件通过51个，另外8个冻结控制器强扰动条件仅通过2个；四步参考用时29秒、前进约0.210米，属于准静态步行验证。

报告按七个技术篇章、二十八个阶段组织，保留各阶段的模型假设、指标定义、参数对照、失败边界和代码入口。实测硬件结果仅限明确标注的IMU与Jetson采集；机械臂接触、平衡及行走结果均来自仿真。现阶段形成的是控制方法实现与物理验证链路，不等同于真实机器人部署或硬实时控制系统。

\noindent\textbf{关键词：}机器人控制；状态反馈；导纳控制；倒立摆；全身逆动力学；双足步行；物理仿真。

\paragraph{本次增补}在原28个阶段之后增加官方G1强化学习行走复现与接触、方向和停止复核，共29个阶段。新增12种固定工况累计245秒仿真与13段视频。官方策略未训练或修改；与自编准静态控制使用不同模型，不作公平算法排名。动态迈步已复现，航向偏移与归零后持续踏步仍需单独处理。
